% questions.tex —— 1977 年·江苏南通地区初试卷 · 题目内容片段（无答案）
% 由 tools/generate_projects.py 依来源总账自 1977/1977_数学真题_江苏南通地区初试_试题.md 生成，勿手改（改源 MD 后重跑）


\gkpart{江苏南通地区初试卷}

\begin{gktip}
\textbf{来源说明}：本题面依据固定源 PDF《1977江苏南通地区初试》第 1 页（原书页码 247）逐题转录。源页为试题汇编材料中的南通地区初试（预选）卷，仅题面、无解答；页首标题「十一、南通地区」系汇编章节标题，按原卷逐字保留为本文 H1。原卷大题以中文数字「一、」至「十、」编号，本库转写为 \textbf{1．}至 \textbf{10．}；第 1、2 题下的小题原卷以「1．2．」编号，转写为（1）（2）。原卷无附图，无附加题/参考题；配套答案与解析为本库编者独立推导，\textbf{非官方标准答案}。
\end{gktip}

\par\medskip

\begin{question}
\gkti{1}回答下列问题：

（1） 当 $a$ 为何值时，$2a>a$；当 $a$ 为何值时，$2a=a$；当 $a$ 为何值时，$2a<a$。

（2） 当 $x$ 为何值时，$|x|=x$。

\end{question}
\begin{question}
\gkti{2}写出下列函数的自变量 $x$ 的取值范围（在实数范围内）：

（1） $y=\dfrac{1}{2-x}$；

（2） $y=\sqrt{2x-1}$。

\end{question}
\begin{question}
\gkti{3}计算：

\nopagebreak[4]
\[
\left(\frac{2}{3}\right)^{2}+(-5.8)^{0}-\left(\frac{8}{27}\right)^{-\frac{2}{3}}+\left(\frac{1}{2}\right)^{-2}.
\]

\end{question}
\begin{question}
\gkti{4}$k$ 为何值时，方程 $x^{2}-2kx+4=0$ 有两个相等的实数根？并求出这方程的根。

\end{question}
\begin{question}
\gkti{5}已知 $\sin\alpha=\dfrac{3}{5}$（$90^\circ<\alpha<180^\circ$），求 $\cos\alpha$、$\operatorname{tg}\alpha$、$\operatorname{ctg}\alpha$ 的值。

\end{question}
\begin{question}
\gkti{6}试写出抛物线 $y=1+4x-2x^{2}$ 的顶点坐标、对称轴方程和函数 $y$ 的极值。

\end{question}
\begin{question}
\gkti{7}解方程

\nopagebreak[4]
\[
\frac{\lg 2x}{\lg(x-4)}=2.
\]

\end{question}
\begin{question}
\gkti{8}求 200 以内的能被 6 整除的所有自然数的和。

\end{question}
\begin{question}
\gkti{9}已知正四棱锥底面边长为 $12\ \mathrm{cm}$，高为 $2\sqrt{7}\ \mathrm{cm}$，求它的侧棱。

\end{question}
\begin{question}
\gkti{10}两根钢绳悬吊货物时，两条钢绳与铅垂线的夹角分别为 $\alpha$ 和 $\beta$，货物的重量为 $p$，试用表达式分别写出钢绳所承受的拉力 $F_1$ 和 $F_2$。

\end{question}